% =======================================================================================
% Draft only: the general to-do list. The specific gaps are red TODOs where they occur in the text;
% this list is what has to be done, improved or checked across the paper. Delete before submission.
% =======================================================================================
\clearpage
\section*{\textcolor{red}{Open Items (draft only, delete before submission)}}
\label{app:open-items}
{\color{red}
\paragraph{Runs and data.}
\begin{itemize}
  \item A noise \acs{psd} estimated with the sources, since a louder noise than the flow's own makes it
    confident and wrong (\cref{app:evaluation-sampler}). (The power/log grids and Euler are in, 2026-11-18.)
\item \cref{app:mcmc} now reads the six packed references of the eval catalogue (2026-11-18: settings,
    one seed per injection, $\hat R$ and \acs{ess}, the posterior figure, cost; Fisher in \cref{app:fisher}); the
    reports/mcmc README describes them since 2026-11-19 (its page, mcmc.html, stays the October archive).
  \item The reported model of each problem (2026-11-19, user; \texttt{paper/scripts/e64.py}) is the better of the
    two runs that trained to near their end, the 64-epoch run and its fork cooled from epoch 32, by eval-catalogue
    loss: e32 for narrow $S=2$, 8 and full $S=16$, main for narrow $S=4$ (epoch 62) and full $S=8$; for full $S=4$
    the Student-t continuation 261103-n4-full-B-t44 at its last checkpoint (its loss does not compare). The
    tables name each model and its epoch. t44 is the best full-band model on every parameter, but whether the
    Student-t frame helps the other problems is open (one seed, one run, S=4); no further runs are planned.
  \item One seed per configuration: the run-to-run spread is unmeasured, so differences of a few
    hundredths of a nat between forks are not significant. A second seed of narrow $S=2$ (about 32 A100
    hours) is a follow-up in \cref{sec:followups} (2026-11-19, user), not a run for this paper.
\end{itemize}

\paragraph{Analysis.}
\begin{itemize}
  \item The overlap offset is no single parameter group (\cref{app:results-validation}, 2026-11-18): draws
    trade $f_0$ against phase. Whether a different chart or loss would remove that is open.
  \item The full band misses loud sources above about \qty{10}{\milli\hertz}, mostly because the flow
    refines the frame only 2--4-fold there (14--17-fold below) (\cref{app:results-mcmc}). Framed (2026-11-17): the prior
    produces few such sources for training; the cure is left to a follow-up.
  \item The joint-MCMC references of narrow $S=8$ and full $S=16$ never left their Fisher start for
    faint sources (\cref{app:results-mcmc}). The prior-started chains (2026-11-16) confirm the
    references from $\snr=7$ up, and bracket $\snr=3$--7 without settling it, so rerunning the
    references with prior starts would not help; the appendix now scores $\snr<3$ against the prior and
    reports the bracket at $\snr=3$--7 (not in the main text for now). Full $S=16$'s prior-started
    chains landed 2026-11-18: same picture below $\snr=7$, but at $\snr\ge7$ they are 1.2--1.7 times wider
    than the reference and still widening, so full $S=16$'s loud ratios are upper bounds
    (\cref{app:results-mcmc}). Settling it would need longer prior-started chains (about 40 CPU hours each).

\end{itemize}

\paragraph{Writing.}
\begin{itemize}
  \item Main text restructured 2026-11-20 (no page limit, for arXiv): the waveform and window
    algebra are in \cref{app:forward-model} and the prior's fits in \cref{app:priors};
    \cref{sec:narrow-band} precedes \cref{sec:method} and holds the source counts and the slots
    (\cref{sec:gw-rates}); \cref{sec:method} runs coordinates, flow matching, architecture,
    training, with \cref{fig:network-architecture} in the main text in place of the schematic;
    calibration and cost are read once for all six models
    (\cref{sec:results-calibration,sec:results-cost}). Still to decide: a results sentence in the
    abstract.
  \item TODO(Giorgio) of 2026-11-20: the \acs{sbc} of the \acs{mcmc} is in \cref{tab:results-summary};
    the low-\acs{snr} question (is the \acs{mcmc} overconfident?) is answered in \cref{sec:validation}
    from \cref{app:results-mcmc}, to be confirmed; the \acs{sbc} heatmap over frequency and \acs{snr} is
    still missing (red TODO in \cref{sec:validation}).
  \item \Cref{sec:background-generative}: add references on large-scale multimodal generation beyond
    images; the sentence that asked for them was dropped from the main text until they exist.
  \item Figures: \texttt{fig\_results.py} was recoloured 2026-11-20 (hue for the band, lightness for
    the source count) but \cref{fig:results} is still its previous render; re-render it where the
    run data lives.
  \item Review every \texttt{\% updated 2026-11-09} marker (stale facts fixed in moved text) and
    keep or revert each.
  \item Check the references marked ``from memory, not yet checked'' in \texttt{references.bib}
    (the method references of the training recipe) against arXiv, like the rest.
  \item Before submission: blind mode, anonymised code link, self-citations in the third person,
    and the acknowledgements commented out for review (back at camera-ready).
  \item Consistency pass done 2026-11-17 (notation, sky frames, British
    spelling, band names: \texttt{narrow-B}/\texttt{full-B} are band and network size, XS--XL sizes only);
    redo after the main-text moves. The Notation appendix was removed 2026-11-18: every symbol and
    acronym is now defined where it first appears.
\end{itemize}
}
